\section{Tightness: the limit map}
\label{sec:tightness}

The constants in the certificate have several kinds of boundary.  An
\emph{absolute pin} is witnessed by a weakened certificate that declares an
attacker win although a legal defender line defeats the declared horizon.  A
\emph{relative pin} attains a counting bound or breaks the stated coupling,
clock, kernel, or witness representation.  It leaves open a verifier with
additional state.  A \emph{syntactic pin} is required to make the certificate
grammar and its recurrences defined.  An \emph{arithmetic attainment} reaches
equality for one selected window or local trace but does not supply a complete
certificate in which every other zone term is absent.  OPEN means that no
general pin and no smaller sound replacement is known.  PARTIAL means that a
restricted reduction is proved and broader forms remain open.

The tables give the current value and one reason for every row.  The reason
column states the scope of the boundary.  A value marked relative must not be
reported as a game-wide optimum.  A value marked absolute must retain the
combined premises named in that row.

\begin{table}[p]
  \centering
  \caption{Limit map, obligation and completion rows.}
  \label{tab:tightness-obligations}
  \scriptsize
  \renewcommand{\arraystretch}{1.12}
  \begin{tabular}{@{}p{.055\textwidth}p{.245\textwidth}p{.205\textwidth}p{.405\textwidth}@{}}
    \toprule
    ID & Parameter and value & Status & One-line reason \\
    \midrule
    R1a & Exact live-role band $8(r-1)$ & Relatively pinned &
      A distance-eight relay chain reaches the protected carrier after
      exactly $r$ defender opportunities, so a smaller band breaks L9$'$. \\
    R1b & Uniform live-role band $8(B-1)$ & OPEN &
      The exact-rank chain does not supply a complete certificate with
      $r=B$ and synchronized terminal, legality, and LOSS data. \\
    R2 & Virgin radius $8(E^{\D}-6)$ for $E^{\D}\ge6$ & Fixed-window
      arithmetic attained; full union OPEN &
      A relay of $E^{\D}-6$ steps attains the radius for one window, but one
      of the seed's other incident virgin windows may still select it. \\
    R3 & Touched guard $\cnt_{\D}(W)+E^{\D}(W)\ge6$ & Fixed-window equality
      attained; full pin OPEN &
      Two existing stones followed by four real fills attain equality, but
      no complete recurrence-labelled certificate excludes every other
      mandatory zone. \\
    \bottomrule
  \end{tabular}
\end{table}

\begin{table}[p]
  \centering
  \caption{Limit map, LOSS and kernel rows.}
  \label{tab:tightness-loss-kernel}
  \scriptsize
  \renewcommand{\arraystretch}{1.12}
  \begin{tabular}{@{}p{.055\textwidth}p{.245\textwidth}p{.205\textwidth}p{.405\textwidth}@{}}
    \toprule
    ID & Parameter and value & Status & One-line reason \\
    \midrule
    R4a & LOSS witness cap at $b=1$: three windows & Relatively pinned &
      The three triangle pairs have transversal number two, while every
      two-window subfamily has a one-cell transversal. \\
    R4b & LOSS witness cap at $b=2$: five windows & Hexo-specific reduction;
      five relatively pinned &
      The five-cycle needs all five edges for transversal number three, and
      hex-window geometry excludes the generic six-edge $K_4$ obstruction. \\
    R5a & Internal T6 scope $\mhs(P)\le b$ & Relatively pinned &
      At transversal number $b+1$ the exact kernel is empty, so it cannot be
      the searched set of a nonempty ordinary internal node. \\
    R5b & Kernel enforcement of $\neg\ownwinnow(P)$ & Absolute for combined
      enforcement only &
      If both the T6 premise and the D9 diagnostic are omitted, the singleton
      kernel can exclude an immediate defender-winning reply. \\
    R5c & Residual test $\tau(F\setminus d)\le b-1$ & Relatively pinned &
      A first hit of two disjoint singleton threats leaves the one remaining
      placement required by the exact minimum-transversal line. \\
    R6 & LOSS contract $\tau(\mathcal T)>b$ plus universal survivor & Absolute
      for the combined contract &
      At equality the defender can hit every named witness; the numeric test
      alone is redundant when the universal survivor clause is still checked. \\
    \bottomrule
  \end{tabular}
\end{table}

\begin{table}[p]
  \centering
  \caption{Limit map, clocks and grammar rows.}
  \label{tab:tightness-clocks}
  \scriptsize
  \renewcommand{\arraystretch}{1.12}
  \begin{tabular}{@{}p{.055\textwidth}p{.245\textwidth}p{.205\textwidth}p{.405\textwidth}@{}}
    \toprule
    ID & Parameter and value & Status & One-line reason \\
    \midrule
    R7 & D17 transition charge: $1+$ child rank and exposure & Relatively
      pinned &
      The current real-only reply can be the missing legality relay or the
      missing window fill, so child-only labels omit one causal event. \\
    R8 & D14 budget: AND $1+\max$, LOSS $b$ & Exact as defined; uniform use
      relative &
      A maximizing finite path uses the current edge and every LOSS remainder
      placement, although indexed guards may use smaller quantities. \\
    R9 & Legality-chain coefficient eight & Relatively pinned &
      The rules legalize a relay at distance exactly eight, and the protected
      occupation and virgin traces attain that step length. \\
    R10 & D15: AND $+1$, OR $+0$, deadline zero, role maximum, protection
      through check & Mixed; leaf-entry and OR-COMPLETION clauses absolute &
      Rank and maximum clauses are limits of the union coupling; dropping
      either named deadline protection admits a concrete false certificate. \\
    R11 & LOSS deadline $p_{\rm leaf}+b+2$ & Absolutely pinned &
      After $b$ hits, a surviving count-four window has two empties and
      completes only on the second following attacker placement. \\
    R12 & D16: AND $1+\max$, LOSS $b$, stopping value zero & Exact as defined;
      uniform use relative &
      The recurrence is the attained path maximum before resolution or first
      attacker entry, while a branch-indexed hazard may be smaller. \\
    \bottomrule
  \end{tabular}
\end{table}

\begin{table}[p]
  \centering
  \caption{Limit map, static coverage and forced-hit rows.}
  \label{tab:tightness-static}
  \scriptsize
  \renewcommand{\arraystretch}{1.12}
  \begin{tabular}{@{}p{.055\textwidth}p{.245\textwidth}p{.205\textwidth}p{.405\textwidth}@{}}
    \toprule
    ID & Parameter and value & Status & One-line reason \\
    \midrule
    R13 & T5 static cover: $B\le3$ with $r_3$ & Local arithmetic attained;
      full pin OPEN &
      Radius three is attained at $B=3$ and the local formula fails at $B=4$,
      but the coordinate trace is not a complete D9 certificate. \\
    R14 & L10 cutoff: first three future threat-creating placements &
      Relatively pinned &
      An explicit fourth placement creates a count-four virgin target while
      lying outside both the initial touched-window set and $r_3$. \\
    R15 & Independently nonempty AND set: at least one legal reply &
      Relative and syntactic &
      A zero-child AND supplies no ghost filler and leaves D14/D16 maxima
      undefined, but deleting nonemptiness alone gives no false-win witness. \\
    R16 & Forced-hit debit: exact-copy $f$ and $Q^{\D}$; full $B,r,E^{\D}$
      retained & PARTIALLY RESOLVED &
      Tight protected gates debit copied hits for named hazards; scalar,
      substituted-hit, slack-pressure, and net-shrinkage claims remain open. \\
    \bottomrule
  \end{tabular}
\end{table}

\subsection{Obligation and completion boundaries}

Rows R1a and R9 describe one mechanism.  A real-only legality chain can
advance by eight at each ordinary defender opportunity.  The coordinate
family attains $8(r-1)$ before a rank-$r$ carrier is checked.  Lowering the
step or exact-rank band breaks the first-protected-occupation argument.  The
uniform replacement $8(B-1)$ is different.  Its soundness follows from
$r\le B$, but no complete terminal certificate realizes all required data
with $r=B$.  Its sharpness remains open.

Rows R2 and R3 distinguish one selected window from the verifier's full
union.  A virgin-window trace uses $E-6$ distance-eight relays and six fills,
attaining $8(E-6)$.  An incident virgin window can still select the same seed
at distance zero, so this does not pin $Z_{\rm virgin}$.  A touched window
with two common defender stones and four real fills attains $2+4=6$.  This
pins the local comparison but lacks the recurrence labels, terminal data, and
exclusion of other zone terms needed for a full countercertificate.

\subsection{LOSS and kernel boundaries}

Every threat window has one or two empties.  At defender budget one, the
triangle family
\[
  \{a,b\},\qquad \{a,c\},\qquad \{b,c\}
\]
has transversal number two and needs all three members.  At budget two, the
five-cycle has transversal number three and deleting any edge drops it to
two.  A six-member minimal obstruction would be $K_4$.  It cannot arise from
Hexo windows: every empty pair is axis-collinear, and four such cells either
violate the three-axis incidence geometry or lie on one line, where an
intervening empty prevents a required two-cell mask.  The three/five caps are
exact for the named-window format, not for every possible witness language.

Rows R5a and R5c are limits of the exact-kernel grammar.  The scope
$\mhs(P)\le b$ keeps the kernel nonempty, and $b-1$ is the placement capacity
after the current reply.  R5b is absolute only under combined enforcement:
if both the T6 premise and D9 diagnostic omit the defender non-win test, a
singleton kernel can exclude an immediate defender completion.  Keeping
either test rejects that certificate.

R6 has the same combined-clause scope.  At
$\tau(\mathcal T)=b$, the defender can hit every named witness.  This defeats
a contract that accepts equality.  If the universal survivor clause is still
checked directly, removing the redundant numeric test has no effect.

\subsection{Clock, deadline, and grammar boundaries}

The R8 and R12 scalar recurrences are attained path maxima.  An AND edge adds
one, a LOSS remainder may use all $b$ placements, and an OR edge adds none.
Window exposure stops on attacker entry.  Indexed clocks may still be smaller
for one role or window.  The D17 transition charge in R7 is needed because the
real reply can legalize the next relay or add a window fill before the child
clock starts.  Another envelope may record that event differently, but it
cannot omit the event.

R10 mixes relative rank clauses with absolute deadline clauses.  AND $+1$,
OR $+0$, and the maximum over roles are limits of the parent-wide union
coupling.  Dropping protection before leaf entry lets the real defender kill
a LOSS witness that remains alive in the ghost.  Omitting an OR-COMPLETION
role lets the real defender occupy the designated winning cell.  Both admit a
false resolution.

R11 is absolute for a separate reason.  Among $b+1$ disjoint count-four
windows, the defender kills $b$ and leaves one with two empties.  The second
following attacker placement completes it, attaining leaf-ply plus $b+2$.
R15 is only syntactic: a zero-child AND supplies no coupling filler and leaves
child maxima undefined, but deleting nonemptiness alone gives no false-win
witness.

\subsection{Static coverage and the forced-hit boundary}

R13 attains radius three at budget three and gives local failure at budget
four, where a touched empty at distance four satisfies $2+4=6$.  Neither
trace supplies the full certificate and clock data for an absolute T5 pin.
R14 is relative for a different reason.  An explicit fourth future attacker
placement creates a threat supported only by future stones while the cell was
initially outside the touched-window set and radius three.

At a protected tight gate, an exact copied hit adds no real-only defender
stone.  The $f$ clock can omit that opportunity, while $Q^{\D}$ charges it
only for a selected window containing the hit and charges full escape and
LOSS remainders.  The scalar $B$, full $r$, and full $E^{\D}$ remain.
A D17 substitute may create a frontier, checkpoint roles may enlarge the
protected set, and slack pressure leaves a quiet placement.  R16 therefore
proves hazard debits, not scalar reduction, free substitution, automatic net
shrinkage, or a slack-pressure rule.

The limit map has three practical readings.  The absolute rows are verifier
safety boundaries under their stated combined premises.  The relative rows
identify where this certificate representation uses all of its available
count.  The OPEN and PARTIAL rows identify changes that require a new proof,
not a smaller constant inserted into the current formulas.
